\documentclass[aspectratio=169,11pt]{beamer}
\usepackage{lecturevideo}   % theme + \say \wait \lvanim \attemptq + \lvwrite \lvzoom ... (visual-style.md)
\title{Lecture title in sentence case}
\subtitle{Course code \textperiodcentered\ Lecture N}
\author{}\institute{}\date{}

% Figure macros live in the preamble (a # inside a frame body breaks beamer).
% Edge-dislocation lattice: #1 = core x. Upper atoms shift by a sigmoid around the core.
\newcommand{\lattice}[1]{%
  \foreach \r in {0,...,7}{\foreach \c in {0,...,15}{
    \pgfmathsetmacro\xx{\c + (\r>3 ? 1/(1+exp(max(-9,min(9,-(#1-\c)*0.9)))) : 0)}
    \pgfmathsetmacro\yy{\r + (\r>3 ? 0.12 : -0.12)*exp(-((\xx-#1)^2)/3)}
    \node[atom,minimum size=7pt] at (\xx,\yy) {};}}
  \draw[accent,dashed] (-0.6,3.5)--(16.6,3.5);
  \node[note,anchor=west] at (16.7,3.5) {slip plane};
  \ifdim#1pt<15.8pt\ifdim#1pt>-0.8pt\node[accent,font=\bfseries\Large] at (#1,3.62) {$\perp$};\fi\fi}

\begin{document}

% Title: \titlepage inside a normal frame so it can carry narration.
\begin{frame}[plain]
\say{One or two sentences: what this lecture covers and why it matters.}
\titlepage
\end{frame}

% Build-up frame: one \say per overlay page, overlay specs must match.
\begin{frame}{A concept, revealed in steps}
\say<1>{First idea.}
\say<2>{Second idea, while the diagram gains a new element.}
\begin{itemize}
\item<1-> first point
\item<2-> second point
\end{itemize}
\end{frame}

% Teaching visuals (references/visual-style.md): the picture first, one colour per idea, a label
% hand-written as it is said (\lvwrite: one ranged \say covers its pages), a halo on what "this"
% is, and a camera zoom (\lvzoom) that pulls back on the next page.
\begin{frame}{Teaching visuals: write, point, zoom}
\say<1>{Picture first: a beam on a pivot, a blue load on the left and a green push on the right.}
\say<2-7>{The blue block is the load, the weight we want to lift.}
\say<8>{Now look closely at the pivot. This point is the fulcrum, and every turning effect is measured from here.}
\say<9>{Pull back, and the balance reads: load times its distance equals effort times its distance.}
\lvzoom<8>{0.5}{0.55}{2}
\centering
\begin{tikzpicture}
  \useasboundingbox (-3.3,-2.1) rectangle (3.3,1.7);  % fixed size: later labels don't shift it
  \draw[ink,line width=3pt] (-3,0) -- (3,0);
  \fill[mute] (0,0) -- (-0.3,-0.5) -- (0.3,-0.5) -- cycle;
  \fill[blu!60] (-2.8,0.05) rectangle ++(0.8,0.8);
  \draw[-{Stealth},grn,very thick] (2.6,1.3) -- (2.6,0.1);
  \node[blu] at (-2.4,1.3) {\lvwrite{2}{6}{load}};
  \only<8->{\lvhalo{(0,-0.05)}\node[note,below=8pt] at (0,-0.5) {fulcrum};}
  \only<9>{\node at (0,-1.7) {$\textcolor{blu}{W\,d_1}=\textcolor{grn}{F\,d_2}$};}
\end{tikzpicture}
\end{frame}

% Derivation frame: physical model on the left, every algebra step on the right with its
% reason, then the physical meaning. Explain each equation mathematically, then intuitively.
\begin{frame}{Deriving the lever rule}
\say<1>{Here's the picture to keep in mind. At this temperature our alloy has split into two phases: some liquid with composition c prime, and some solid with composition c double prime. The overall composition c sits somewhere in between.}
\say<2>{Step one is just bookkeeping. All the B atoms have to be somewhere, so the total amount of B equals the B in the liquid plus the B in the solid. Each phase contributes its fraction times its composition.}
\say<3>{Step two: there are only two phases, so their fractions have to add up to one. That lets us write the liquid fraction as one minus the solid fraction.}
\say<4>{Now substitute that into step one, so the only unknown left is the solid fraction.}
\say<5>{Collect the terms and divide. The solid fraction is c minus c prime, over c double prime minus c prime.}
\say<6>{Here the solid contains more B than the liquid. At this temperature, both phase compositions are fixed. An alloy with more B must therefore contain more solid. Check the liquid endpoint: when c equals c prime, the solid fraction is zero.}
\say<7>{The tie line gives the phase compositions. Mass balance converts those compositions into phase amounts. Explain why composition and phase fraction are different quantities. Pause longer if you need to.}
\wait<8>{6} % Keep the completed diagram and equations visible while the learner connects them.
\begin{columns}[T]
\begin{column}{0.42\textwidth}\centering
\begin{tikzpicture}[x=0.5cm,y=0.5cm]
  \draw[tie] (0,2)--(8,2);
  \fill[blu] (0,2) circle (3pt) node[lab,above=3pt] {$c'$} node[note,below=4pt] {liquid};
  \fill[grn] (8,2) circle (3pt) node[lab,above=3pt] {$c''$} node[note,below=4pt] {solid};
  \fill[accent] (3,2) circle (3pt) node[lab,above=3pt,accent] {$c$};
  \only<6->{\draw[accent,thick] (3,1.2)--(2.6,0.6)--(3.4,0.6)--cycle;
    \draw[accent,<->] (0,-0.3)--node[note,below] {arm for $X_S$} (3,-0.3);
    \draw[mute,<->] (3,-1.6)--node[note,below] {arm for $X_L$} (8,-1.6);}
\end{tikzpicture}
\end{column}
\begin{column}{0.56\textwidth}\small
\begin{tabular}{@{}l@{\quad}l@{}}
\uncover<2->{$c = X_L\,c' + X_S\,c''$} & \uncover<2->{\color{mute}\footnotesize mass balance on B}\\[3pt]
\uncover<3->{$X_L = 1 - X_S$} & \uncover<3->{\color{mute}\footnotesize only two phases}\\[3pt]
\uncover<4->{$c = (1-X_S)\,c' + X_S\,c''$} & \uncover<4->{\color{mute}\footnotesize substitute}\\[3pt]
\uncover<5->{$\boxed{X_S = \dfrac{c-c'}{c''-c'}}$} & \uncover<5->{\color{mute}\footnotesize collect $X_S$, divide}
\end{tabular}\par\medskip
\uncover<6->{{\color{accent}Physically:} fixed phase compositions; more overall B requires
more of the B-rich solid. Check: $c=c'$ gives $X_S=0$.}
\end{column}
\end{columns}
\end{frame}

% Animation: \lvanim{first}{frames}{code with \lvt from 0 to 1}, narrated by ONE \say spanning it.
% Example: edge dislocation glide. Upper half-crystal atoms shift by a sigmoid around the core.
\begin{frame}{Animation: an edge dislocation glides}
\say<1>{Here is a perfect crystal. A dislocation is about to enter from the left, along the slip plane.}
\say<2-13>{Under shear, the core glides to the right. Each atom shifts only slightly as it passes, yet once the dislocation exits, the top half has slipped by one Burgers vector, leaving a step on the surface.}
\say<14>{The dislocation has left the crystal. The top half remains displaced by one Burgers vector. Local motion of the defect has produced permanent slip. Explain what moved locally and what changed across the crystal. Pause longer if you need to.}
\wait<15>{6}
\centering
\begin{tikzpicture}[x=0.55cm,y=0.55cm]

\only<1>{\lattice{-4}}
\lvanim{2}{12}{\pgfmathsetmacro\cx{-4+23*\lvt}\lattice{\cx}}
\only<14->{\lattice{19}} % Hold the final animation state during consolidation.
\end{tikzpicture}
\end{frame}

% DEMONSTRATION ONLY, not sourced homework. Before generation, ask for the actual problem
% set and replace this item with a selected original question, figure, and exercise ID.
% The hold is a buffer for manually pausing the player; it is not a 15-second solve limit.
\attemptq[15]{Example only: use mass balance}
  {At fixed $T_3$, liquid has $c'=20$ wt.\% B and solid has $c''=85$ wt.\% B.
   The solid mass fraction is $X_S=0.30$.\par\smallskip
   (a) Find the overall composition. (b) Check a claim that the liquid fraction is also 0.30.
   Justify your conclusion using mass balance.}
  {At a fixed temperature T three, the liquid contains twenty weight percent B and the solid
   contains eighty five. The solid mass fraction is zero point three. Find the overall
   composition. Then check the claim that the liquid fraction is also zero point three.
   Justify your conclusion using mass balance. Allow about three minutes for your attempt.}
  {$X_L=1-0.30=0.70$.\par\medskip
   $c=0.70(20)+0.30(85)=39.5$ wt.\% B.\par\medskip
   The proposed fractions sum to $0.60$, not $1$, so they cannot describe these two phases.
   If the fractions were swapped, $c=0.30(20)+0.70(85)=65.5$ wt.\% B, a different alloy.}
  {Only two phases are present, so their mass fractions add to one. The liquid fraction is
   therefore zero point seven. Multiply each mass fraction by its phase composition and add:
   the overall composition is thirty nine point five weight percent B. Assigning zero point
   three to both phases accounts for only sixty percent of the mass. Swapping their fractions
   would instead give sixty five point five weight percent B.}

% Also deliver recall.md and recall-answers.md; see references/learning-design.md.
\end{document}
